% Generated by scripts/build_reports.py; edit reviewed Markdown inputs.
% Compile with XeLaTeX.
% Options for packages loaded elsewhere
\PassOptionsToPackage{unicode}{hyperref}
\PassOptionsToPackage{hyphens}{url}
\PassOptionsToPackage{dvipsnames,svgnames,x11names}{xcolor}
\documentclass[
  11pt,
]{article}
\usepackage{xcolor}
\usepackage[margin=1in]{geometry}
\usepackage{amsmath,amssymb}
\setcounter{secnumdepth}{-\maxdimen} % remove section numbering
\usepackage{iftex}
\ifPDFTeX
  \usepackage[T1]{fontenc}
  \usepackage[utf8]{inputenc}
  \usepackage{textcomp} % provide euro and other symbols
\else % if luatex or xetex
  \usepackage{unicode-math} % this also loads fontspec
  \defaultfontfeatures{Scale=MatchLowercase}
  \defaultfontfeatures[\rmfamily]{Ligatures=TeX,Scale=1}
\fi
\usepackage{lmodern}
\ifPDFTeX\else
  % xetex/luatex font selection
\fi
% Use upquote if available, for straight quotes in verbatim environments
\IfFileExists{upquote.sty}{\usepackage{upquote}}{}
\IfFileExists{microtype.sty}{% use microtype if available
  \usepackage[]{microtype}
  \UseMicrotypeSet[protrusion]{basicmath} % disable protrusion for tt fonts
}{}
\makeatletter
\@ifundefined{KOMAClassName}{% if non-KOMA class
  \IfFileExists{parskip.sty}{%
    \usepackage{parskip}
  }{% else
    \setlength{\parindent}{0pt}
    \setlength{\parskip}{6pt plus 2pt minus 1pt}}
}{% if KOMA class
  \KOMAoptions{parskip=half}}
\makeatother
\usepackage{longtable,booktabs,array}
\usepackage{caption}
\captionsetup[table]{skip=6pt}
\newcounter{none} % for unnumbered tables
\usepackage{calc} % for calculating minipage widths
% Correct order of tables after \paragraph or \subparagraph
\usepackage{etoolbox}
\makeatletter
\patchcmd\longtable{\par}{\if@noskipsec\mbox{}\fi\par}{}{}
\makeatother
% Allow footnotes in longtable head/foot
\IfFileExists{footnotehyper.sty}{\usepackage{footnotehyper}}{\usepackage{footnote}}
\makesavenoteenv{longtable}
\setlength{\emergencystretch}{3em} % prevent overfull lines
\providecommand{\tightlist}{%
  \setlength{\itemsep}{0pt}\setlength{\parskip}{0pt}}
\usepackage{amsmath,amssymb}
\usepackage{xurl}
\usepackage{fancyhdr}
\pagestyle{fancy}
\fancyhf{}
\fancyhead[L]{\small Pure Thought | Revised research specifications}
\fancyfoot[L]{\small 2026-09-17}
\fancyfoot[R]{\small\thepage}
\setlength{\headheight}{15pt}
\setlength{\emergencystretch}{3em}
\allowdisplaybreaks
\usepackage{bookmark}
\IfFileExists{xurl.sty}{\usepackage{xurl}}{} % add URL line breaks if available
\urlstyle{same}
% fallback for those not using the hyperref driver hyperxmp:
\makeatletter
\@ifundefined{xmpquote}{\newcommand{\xmpquote}[1]{#1}}{}
\makeatother
\hypersetup{
  pdftitle={24. Topological code verification and noise-specific decoder evaluation},
  pdfauthor={Pure Thought research portfolio},
  colorlinks=true,
  linkcolor={blue},
  filecolor={Maroon},
  citecolor={Blue},
  urlcolor={blue},
  pdfcreator={LaTeX via pandoc}}

\title{24. Topological code verification and noise-specific decoder
evaluation}
\author{Pure Thought research portfolio}
\date{Revised 17 September 2026}

\begin{document}
\maketitle
\thispagestyle{empty}
\clearpage

\section{24. Topological code verification and noise-specific decoder
evaluation}

\textbf{Revised specification --- 2026-09-11.}
\href{https://github.com/rick-stevens-ai/Pure-Thought/blob/main/reviews/24-Topological-Quantum-Error-Correction.md}{Detailed
critique} \ensuremath{\cdot}
\href{https://github.com/rick-stevens-ai/Pure-Thought/blob/main/archive/original-PRDs/24-Topological-Quantum-Error-Correction.md}{Original
description} \ensuremath{\cdot}
\href{https://github.com/rick-stevens-ai/Pure-Thought/blob/main/reviews/EVIDENCE-STANDARD.md}{Evidence
standard}

\subsection{Plain-language guide}

\subsubsection{The problem in everyday terms}

Imagine storing information in loops spread across a lattice so that a
small local disturbance usually cannot change the stored value. A
topological quantum code uses a quantum version of this idea. Local
measurements reveal the endpoints of error patterns, and a decoder tries
to join those endpoints correctly. The subtlety is that two repairs can
have the same endpoints yet differ by a loop that changes the logical
information.

\subsubsection{Key terms}

\begin{itemize}
\item
  \textbf{Toric code} --- A quantum error-correcting code on a lattice
  with periodic boundaries, topologically equivalent to a torus.
\item
  \textbf{Surface code} --- A related geometric code with carefully
  chosen open boundaries.
\item
  \textbf{Logical operator} --- An operation changing the encoded
  information while remaining consistent with the code's checks.
\item
  \textbf{Noncontractible loop} --- A loop winding around the torus that
  cannot be shrunk to a point within it.
\item
  \textbf{Homology} --- A mathematical way of distinguishing cycles
  modulo boundaries; here it identifies whether a closed error pattern
  acts logically.
\item
  \textbf{Syndrome} --- Check outcomes locating the boundary of an error
  chain.
\item
  \textbf{Recovery} --- An operation chosen to cancel the error;
  matching the syndrome is necessary but does not ensure logical
  success.
\item
  \textbf{Minimum-weight matching} --- A graph optimization method used
  in certain decoders to pair defects efficiently.
\item
  \textbf{Logical failure probability} --- The probability that error
  and recovery together change encoded information.
\item
  \textbf{Threshold} --- An asymptotic noise boundary for successful
  suppression of logical errors in a specified code-and-decoder family.
\end{itemize}

\subsubsection{Why this matters}

Geometric codes connect an intuitive picture of local errors to precise
algebraic guarantees. They are important models for protecting quantum
information using structured measurements. Understanding logical
failure, rather than only syndrome cancellation, is essential for
assessing whether increasing the code size actually improves
reliability.

\subsubsection{What progress would mean}

A useful result would prove the code's parameters and measure a fixed
decoder under a declared noise model. An estimated threshold would
remain specific to that model and decoder, not a universal device error
tolerance.

\subsection{Core question}

Can a toric-code family's algebraic parameters be proved, and can a
fixed decoder's logical failure probability be evaluated under a
precisely specified error model?

\subsection{Scope and assumptions}

Use an L\ensuremath{\times}L square torus with qubits on edges and
L\ensuremath{\ge}3. Start with independent Z errors of probability p and
perfect syndrome extraction. Specify the decoder, tie-breaking rule, and
whether recovery weights are uniform or likelihood based.

Keep exact code properties separate from statistical decoder
performance. Planar boundaries, color codes, measurement noise, and
circuit-level noise are distinct extensions.

\subsection{Mathematical target}

Over F\ensuremath{{}_{2}} use
C\ensuremath{{}_{2}}\ensuremath{\to}\ensuremath{\partial}\ensuremath{{}_{2}}
C\ensuremath{{}_{1}}\ensuremath{\to}\ensuremath{\partial}\ensuremath{{}_{1}}
C\ensuremath{{}_{0}} with
\ensuremath{\partial}\ensuremath{{}_{1}}\ensuremath{\partial}\ensuremath{{}_{2}}=0,
H\_X=\ensuremath{\partial}\ensuremath{{}_{1}},
H\_Z=\ensuremath{\partial}\ensuremath{{}_{2}}\ensuremath{{}^{T}}. Then
n=2L\ensuremath{{}^{2}}, k=2, and prove d=L using noncontractible primal
and dual cycles.

For Z error e and recovery r, success means e+r\ensuremath{\in}im
\ensuremath{\partial}\ensuremath{{}_{2}}, not merely
\ensuremath{\partial}\ensuremath{{}_{1}}(e+r)=0. A threshold assertion
concerns the limit of P\_fail(L,p) as
L\ensuremath{\to}\ensuremath{\infty} for this fixed noise/decoder
family; a finite crossing is an estimate.

\subsection{Required outputs and proof obligations}

Export lattice incidence data, ranks and logical bases, distance
argument, decoder specification, and failure counts with confidence
intervals. Exhaustive small-size evaluation can give an exact failure
polynomial in p.~Monte Carlo output must retain trials and failure
counts, including zero-failure upper intervals.

\subsection{Validation and rejection controls}

Check n,k,d on several small tori. Verify logical anticommutation via
primal/dual intersection. Enumerate all low-weight errors on tiny
instances and test the stabilizer-membership success criterion. Validate
a matching solution's optimality separately from logical decoding
success.

\subsection{Milestones and research extension}

Baseline: exact construction and distance proof with a tested decoder.
Strong: reproducible finite-size failure curves and uncertainty-aware
threshold estimates. Research extension: a rigorous threshold bound or
noisy-measurement/circuit model, with no transfer of numerical
thresholds across noise conventions.

Milestones are acceptance gates, not calendar promises. Before a run,
freeze the concrete instance/domain, tolerance, compute budget, and
permitted dependencies. A bounded search that fails to find a result
must preserve an unresolved outcome.

\subsection{Starting instruction}

Work on the baseline above. First make every input and convention
explicit, then validate the reference case. Build the checker alongside
the calculation. Label numerical evidence, finite-model certificates,
and continuum theorems separately; return the strongest justified result
and identify the exact unmet proof obligation.

\subsection{Primary references and their role}

\begin{itemize}
\tightlist
\item
  \href{https://arxiv.org/abs/quant-ph/0110143}{Dennis, Kitaev, Landahl
  and Preskill, Topological quantum memory} --- Homological coding and
  model-dependent error-correction thresholds.
\end{itemize}

References support definitions or limitations; the scope choices and
proposed milestones are editorial recommendations, not claims that these
results have already been obtained.

\subsection{Displayed formulation and claim boundary}

\href{https://github.com/rick-stevens-ai/Pure-Thought/blob/main/PRDs/24-Topological-Quantum-Error-Correction.md}{Canonical
specification}.

\[C_2\xrightarrow{\partial_2}C_1\xrightarrow{\partial_1}C_0,
\quad\partial_1\partial_2=0,\quad H_X=\partial_1,\quad H_Z=\partial_2^{\mathsf T},\]

\[n=2L^2,\quad k=2,\quad d=L\quad(L\ge3),\qquad
e+r\in\operatorname{im}\partial_2\quad\text{for successful }Z\text{-error recovery}.\]

A vanishing residual syndrome only gives \(e+r\in\ker\partial_1\) and
can hide a nontrivial logical cycle. Statistical decoder curves require
the fixed noise model and decoder; finite-size crossings are threshold
estimates, not limit theorems.

\subsection{Detailed review and rationale}

\href{https://github.com/rick-stevens-ai/Pure-Thought/blob/main/PRDs/24-Topological-Quantum-Error-Correction.md}{Revised
problem} \ensuremath{\cdot}
\href{https://github.com/rick-stevens-ai/Pure-Thought/blob/main/archive/original-PRDs/24-Topological-Quantum-Error-Correction.md}{Preserved
original} \ensuremath{\cdot}
\href{https://github.com/rick-stevens-ai/Pure-Thought/blob/f854d4e12b55c1f5b5c2a964b75f2d23496d981e/PRDs/24-Topological-Quantum-Error-Correction.md}{Pinned
upstream source}

Line references below refer to the preserved original at commit
\texttt{f854d4e12b55}, not the rewritten file. ``Blocking'' means the
original issue can invalidate the scientific target or claimed
certificate; ``major'' means it materially changes scope, inference, or
acceptance criteria. These are critiques of specifications, not results
of running the embedded research-code sketches.

\subsubsection{Detailed findings}

\begin{enumerate}
\def\labelenumi{\arabic{enumi}.}
\item
  \textbf{Blocking --- noise mismatch (17, 88, 1003).} A
  \textasciitilde10.9\% figure is not a universal depolarizing/circuit
  threshold. Noise channel, measurement assumptions, and decoder
  determine the comparison.
\item
  \textbf{Major --- coefficient field (97--107).} CSS homology is
  computed over F\ensuremath{{}_{2}}; writing
  H\ensuremath{{}_{1}}=Z\ensuremath{{}^{2}} without identifying the
  coefficient change confuses topological and binary-code data.
\item
  \textbf{Blocking --- syndrome versus correction (110--113, 133--136).}
  Two chains with the same boundary may differ by a logical loop.
  Success requires trivial homology of error plus recovery.
\item
  \textbf{Major --- certification boundary (49--52, 136).} Monte Carlo
  threshold fits are not algebraic proofs, and prohibiting numerical
  optimization conflicts with matching and simulation plans. The
  revision allows useful numerical work while reserving proof language
  for justified conclusions.
\end{enumerate}

\subsubsection{What the revision preserves and changes}

The central motivation is retained, but the executable question is now:
\textbf{Can a toric-code family's algebraic parameters be proved, and
can a fixed decoder's logical failure probability be evaluated under a
precisely specified error model?} The revised scope chooses a concrete
starting model; those choices are proposed research-design decisions,
not facts inferred from the original. The baseline can succeed without
resolving the full open research program.

The replacement also removes the original implementation sketches from
the active specification. They remain in the archive for comparison, but
have not been made executable or certified. New implementations should
derive their equations from the revised definitions rather than copy
unchecked prototypes.

\subsubsection{Acceptance criteria to use instead}

Baseline: exact construction and distance proof with a tested decoder.
Strong: reproducible finite-size failure curves and uncertainty-aware
threshold estimates. Research extension: a rigorous threshold bound or
noisy-measurement/circuit model, with no transfer of numerical
thresholds across noise conventions.

\subsubsection{Sources supporting the corrections}

\begin{itemize}
\tightlist
\item
  \href{https://arxiv.org/abs/quant-ph/0110143}{Dennis, Kitaev, Landahl
  and Preskill, Topological quantum memory} --- Homological coding and
  model-dependent error-correction thresholds.
\end{itemize}

\section{Evidence and execution standard}

This standard applies to the revised specifications. It is a proposed
contract for future research runs; this editorial revision does not
supply the scientific certificates themselves.

\subsection{What ``pure thought'' means here}

Use explicitly stated mathematical models and deductions without fitting
conclusions to experimental outcomes. Synthetic examples, exact model
inputs, analytic references, numerical search, and established software
are allowed. Physical constants, electronic-structure approximations,
constitutive laws, and empirical parameter sets must be identified when
used. A model-derived result is conditional on that model.

The original ban on legacy software conflicts with its proposed reliance
on NumPy, SciPy, CVXPY, nauty, and other libraries. Reimplementing every
optimizer and algebra system is not necessary for scientific
independence. Prefer a small independently checkable certificate and a
documented trusted base. A fresh implementation can be a project
objective, but must not be confused with a mathematical proof or novelty
claim.

\subsection{Four evidence levels}

{\def\LTcaptype{none} % do not increment counter
\begin{longtable}[]{@{}
  >{\raggedright\arraybackslash}p{(\linewidth - 4\tabcolsep) * \real{0.3333}}
  >{\raggedright\arraybackslash}p{(\linewidth - 4\tabcolsep) * \real{0.3333}}
  >{\raggedright\arraybackslash}p{(\linewidth - 4\tabcolsep) * \real{0.3333}}@{}}
\toprule\noalign{}
\begin{minipage}[b]{\linewidth}\raggedright
Label
\end{minipage} & \begin{minipage}[b]{\linewidth}\raggedright
What it establishes
\end{minipage} & \begin{minipage}[b]{\linewidth}\raggedright
What it does not establish
\end{minipage} \\
\midrule\noalign{}
\endhead
\bottomrule\noalign{}
\endlastfoot
Numerical evidence & A reproducible computation at stated precision,
sample size, and discretization & Exact truth, continuum coverage, or a
universal theorem \\
Finite-model certificate & An exact/validated statement about the
specified finite model or finite relaxation & An omitted tail,
thermodynamic limit, or physical realization \\
Conditional mathematical theorem & A proof over the stated domain with
every hypothesis and error controlled & Truth outside those assumptions
or a demonstrated physical system \\
Formal proof & A checked theorem in a named proof assistant, with
imported axioms/dependencies listed & Correct physical modeling merely
by virtue of formalization \\
\end{longtable}
}

Formalization is an optional extra deliverable unless the specific
project makes it mandatory. High precision, JSON output, solver status
and small residuals do not by themselves confer any of these proof
levels.

\subsection{Claim directions}

For minimization over a larger feasible set, the optimum is a lower
bound on the original minimum. A physical feasible point supplies an
upper bound. For maximization the directions reverse. Track feasibility
of the primal/dual witnesses separately from their objective values.

Infeasibility of an inner ansatz excludes only that ansatz.
Infeasibility of an outer relaxation can exclude the original problem if
the map into the relaxation is proved and all approximations preserve
necessity. Passing necessary tests means ``not excluded''; existence
needs a constructive sufficient argument.

A smallest observed example is not a minimality theorem. A complete
catalog must define the universe, equivalence relation, boundaries, and
proof that pruning misses nothing. A finite-size numerical trend does
not prove an asymptotic limit.

\subsection{Minimum run manifest}

Before calculation, record the problem ID and revision hash, exact
inputs and conventions, quantified domain, equivalence relation, target
observable, accepted evidence level, tolerances, and resource limits.
Record dependencies and precision/rounding mode. Where randomization is
used, preserve seeds, trial counts and raw sufficient statistics.
Resource limits are planning decisions, not evidence that a mathematical
solution cannot exist.

An output bundle should include:

\begin{itemize}
\tightlist
\item
  The statement being claimed, including all assumptions and
  quantifiers.
\item
  Construction or proof object, plus an independently runnable verifier
  and its trust assumptions.
\item
  Numerical/error ledger separating roundoff, discretization,
  truncation, optimization, sampling and modeling error where relevant.
\item
  Positive reference cases and negative controls that the verifier must
  reject.
\item
  Coverage: included domain, omitted domain, unresolved regions and
  termination reason.
\end{itemize}

These are requirements for future implementations, not a promise of
polynomial-size certificates for every problem.

\subsection{Checking specific certificate types}

\textbf{LP/SDP/SOS:} verify equality constraints, objective direction,
and cone feasibility using exact arithmetic or outward-rounded
intervals. Rationalizing a numerical solution can break PSD constraints.
Correct residuals with justified norm bounds; do not silently project
and reuse the old objective. Distinguish a proof of a polynomial
approximation from a proof for the original function.

\textbf{Interval numerics:} state function enclosures, domains and
outward rounding. A small root box needs existence and, if claimed,
uniqueness. Cover the complement for a global enumeration. Include
singularities and tails rather than discarding them numerically.

\textbf{SAT/SMT:} a solver's UNSAT label is not a certificate. Preserve
a proof supported by a checker and the proof that the encoding, symmetry
reduction and blocked solutions match the mathematical universe. DRAT
applies to an appropriate Boolean CNF proof, not automatically to every
real-arithmetic or SDP claim.

\textbf{Statistics:} specify the estimand, dependence assumptions,
confidence procedure and bias. Zero observed failures does not prove
zero failure probability. MCMC error bars based on estimated
autocorrelation are not generally rigorous mixing guarantees.

\subsection{Milestones and novelty}

Baseline means a working, independently validated result on a concrete
reference problem. Strong means a broader or sharper result with its
claimed evidence level. Research extension means a possible novel result
whose feasibility and significance need review. A negative theorem can
be valuable; an unsuccessful search can also be useful if its scope is
documented, but the two are not interchangeable.

Do not promise ``publication-quality'' in a fixed number of months.
Before claiming novelty, conduct a focused current literature comparison
against the exact theorem/model/metric. The sources in this revision
check key definitions and limitations; they are not an exhaustive
novelty audit across all 30 fields.

\end{document}
